\begin{theorem}[A quadratic controlled witness for two-entry caches]
\label{thm:c2quadraticwitness}
For the adopted algorithm with $c=\ell=2$ and $n\ge3$, put
\[
 T_n=16n^2,\qquad \gamma_n=(4n)^{-T_n}.
\]
For any prescribed labelled canonical state $h=H(V,2)$,
\[
 \inf_x\widehat P_2^{T_n}(x,\{h\})\ge\gamma_n,
 \qquad
 \sup_x\|\widehat P_2^t(x,\cdot)-\pi\|_{\TV}
 \le(1-\gamma_n)^{\lfloor t/T_n\rfloor},
\]
where the infimum permits every weakly connected initial
state with arbitrary finite integer ages.
Consequently
\[
 t_{\rm mix}(\varepsilon)
 \le T_n\left\lceil\gamma_n^{-1}
                  \log(1/\varepsilon)\right\rceil.
\]
The polynomial controlled-path length does not assert a
polynomial mixing time: the probability of the prescribed
path is exponentially small.
\end{theorem}

\begin{proof}
Choose the largest label of the prescribed canonical order
as the owner identity $a$. Write
\[
 I=\{v\ne a:a^0\in C_v\},\qquad k=|I|.
\]
We construct a universal zero-age anchor through at most
$n-1$ strict increases of $k$. Temporary losses during a
macro are permitted; every completed macro strictly
increases its endpoint value.

We first record two elementary services. If $a$ caches
a nonholder target $z$, activate only $a$, retaining $z$
as extra while its partner differs from $z$. At any holder
responder require the full reply, which contains $a$ and,
once learned, $z$. The self identity $a$ is filtered at
its owner, so at most one usable new reply fits the
mandatory vacancy. Thus $z$ survives, the owner's cache
does not grow, first visits establish $z$ at responders,
and repeats shrink the cache. Within $n+2$ events some
nonholder, at latest $z$, receives fresh $a^0$.
All previously zero holders are passive and protected.
If all other $a$ copies were zero at the start, they
remain zero throughout this owner-only service.

Second, suppose $Q$ holds $a^0$ and the owner is frozen.
Activate only $Q$, choosing its residual partner rather
than $a$ in a tie, and send $a$ as extra. Stop before
an exchange to a nonholder if one is encountered.
At every internal holder contact the full reply refreshes
the initiator's anchor to zero, and the passive anchor
stays zero by minimum merging. After deleting the anchor,
these are exactly capacity-one singleton exchanges.
Every visited responder learns $Q$; a repeat returns
$Q$ as self and clears the residual cache.
Within $n$ events $C_Q=\{a^0\}$ or a nonholder is
encountered. The owner's cache stays unchanged.

If a positive-age copy $a^A$ exists at $p$, choose an
actual oldest partner. If it is $a$, one exchange makes
the owner cache the nonholder $p$; the first service gives
a strict $k$ gain within another $n+2$ events.
Otherwise $C_p=\{a^A,u^B\}$ and the request forces
$C_u=\{p^0,a^b\}$. If $u$ was a zero holder, its full
reply gives $p$ a zero and the gain is immediate.
Otherwise $b>0$, so $u$'s next oldest partner is $a$.
After $u\to a$, the owner caches the nonholder $u$,
and the same service gives a gain. This costs at most
$n+4$ events, and no earlier zero holder is lost.
Repeat this rule whenever positive copies remain.
If no $a$ descriptor exists anywhere, weak connectivity
forces a nonempty owner cache; one owner activation
creates a fresh zero. This also counts as a strict gain.

It remains to give an $O(n)$ gain macro when all $a$
descriptors are zero and $k<n-1$. If the owner is
nonempty, inspect an actual oldest partner $Q$.
A nonholder partner gives a gain immediately.
For $Q\in I$, freeze the owner and perform the stopped
residual clearing at $Q$. If it clears, exchange
$a\to Q$. The sole reply $a^0$ is self-filtered,
so the owner loses $Q$ and gains no usable entry.
At $Q$, fresh $a^0$ and at most one extra fit capacity
two. Thus all zeros survive, and the owner loses
one cache entry. Repeat if necessary, at most twice.

If instead the stopped clearing encounters a nonholder
$x$, perform $Q\to x$ and then immediately $a\to Q$.
The first event forces
\[
 C_x=\{Q^0,a^1\}.
\]
Although $Q$ may temporarily lose its zero, the second
event restores it with protected fresh $a^0$.
The owner's choice of $Q$ is still legal because
its cache was frozen during the clearing.
Every other old zero holder was passive and protected.
Now $x\to a$ is forced, since $a^1$ is older than $Q^0$.
This deletes the only new positive copy and makes the
owner cache fresh $x^0$. All old $k$ zero copies are
restored, and all $a$ copies are again zero.
The cached-target owner service adds a zero within
$n+2$ further events. There can have been at most one
earlier owner-head clearing before this branch.
A safe bound for either emptying the owner or gaining
a zero in this phase is therefore $3n+6$ events.

Suppose it ends with the owner empty and no gain.
If a holder caches a nonholder, the four-exchange
gadget of Lemma~\ref{lem:c2infection} gives a gain
with all $a$ copies zero at its endpoint.
Otherwise each holder cache is $a^0$ plus at most
one residual identity from $I$. Weak connectivity,
the empty owner, and the absence of any $a$ descriptor
in the nonholder set
$J=V\setminus(I\cup\{a\})$
supply an incoming boundary descriptor from $J$ to $I$.

Use the incoming-boundary distinction in
Theorem~\ref{thm:thesisc2}, with one modification:
normalize only its selected root $Q$ by the residual
clearing service, rather than making its whole component
a star. This costs at most $n$ events, keeps the outside
boundary initiator frozen, and ends with
$C_Q=\{a^0\}$. The residual weak-component partition
is preserved. If $Q$'s component is nontrivial, its
connectivity and $Q$'s empty residual cache imply
that some other member $v$ still has
$C_v=\{a^0,Q^B\}$ for a finite $B\ge0$.
This is sufficient for Lemma~\ref{lem:c2infection};
that lemma does not require $B=0$.

For completeness, a boundary initiator $P\in J$ whose
oldest partner lies in $J$ first sends the target
$Q\in I$ as extra. That responder is forced to cache
$P^0,Q$, so it can next contact $Q$.
After the root reply $a^0$, it replaces $Q$ as a zero
holder and has a residual edge to nonholder $P$.
The four-exchange gadget therefore gains a zero.
If $P$ instead contacts $Q$ directly, a singleton
$P$ gains immediately. For $C_P=\{Q,X\}$ the exchange
replaces zero holder $Q$ by $P$, which caches
$a^0,X$. If $X\in J$, apply the infection gadget.
If $Q$'s residual component was nontrivial, use the
incoming neighbor $v$ of the now-nonholder $Q$
and apply the same gadget.

The only remaining branch has $Q$ as an isolated
residual component and $X$ in another holder component.
The exchange replaces isolated $Q$ by $P\to X$,
strictly reducing the number of residual components.
It preserves $k$, the empty owner, the all-zero
property of every $a$ copy, and residual closure.
Crucially, this branch costs only one event:
the isolated $Q$ needed no residual clearing.
At most $n-1$ such reductions precede a gain.
The final gain attempt costs at most $n+6$ events,
including root clearing, any preparatory outside
exchange, and the four-event infection.
Thus a safe bound for this incoming phase is $2n+6$.
Together with the owner phase, each all-zero strict-gain
macro costs at most $5n+12$.

The positive-copy and all-zero macros all increase
the same integer $k$. At most $n-1$ such gains produce
a universal zero anchor, with conservative total
budget
\[
 1+(n-1)(5n+12).
\]
The extra initial event is harmless overcounting.

Finally put $D=2n(n-1)$. The capacity-one theorem
normalizes any residual component to a prescribed
zero-age rooted star in at most $D$ events, and the
sum of these budgets across a partition is also at
most $D$. Proposition~\ref{prop:anchorcanonical}
can therefore be implemented with at most
\[
 2(D+1)+D+4(n-2)+D
 \le 4D+4n+2=8n^2-4n+2
\]
events: empty the owner, normalize the residual
components once, merge their roots by the four-event
table, and normalize the connected remainder once.
The root kept during merging stays singleton
$\{a^0\}$ and unprocessed component roots are
untouched, so no intermediate renormalization is
needed. Choose the residual star root as the
second-largest prescribed label. Together with
empty $a$, the result is exactly the desired
$H(V,2)$; no root-change step is needed.

The entire controlled path has length at most
\[
 1+(n-1)(5n+12)+(8n^2-4n+2)
 =13n^2+3n-9<16n^2.
\]
Pad to $T_n$ by activating the canonical empty root.
Every prescribed event has probability at least
$1/(4n)$: the initiator costs $1/n$, an oldest tie
at most a factor two, and a possible passive victim
choice at most a factor two. For $c=\ell=2$ the
request extra and complete reply are forced, and
an active victim, when needed, is unique.
Multiplying these probabilities gives the asserted
common-atom bound. Lemma~\ref{lem:atom} proves the
total-variation estimate, and
$(1-\gamma)^m\le e^{-\gamma m}$ gives the mixing bound.
\end{proof}

